\section{Your First Bare-Metal Rust Program}
\label{sec:bare-metal-program}

% Section 01.02: Your First Bare-Metal Rust Program

After understanding why Rust is compelling for embedded systems, let's create our first bare-metal Rust program. This example will run without an operating system, directly on the microcontroller hardware.

\subsection{Creating a Minimal Embedded Project}

First, create a new directory for our project:

\begin{minted}{bash}
$ mkdir blink_bare_metal
$ cd blink_bare_metal
$ cargo init --name blink_bare_metal
\end{minted}

This creates a standard Rust project structure. However, for embedded development, we need to modify it significantly.

\subsection{Configuring for bare\_metal: \texttt{no\_std}}

Before we begin, let's understand what a bare-metal embedded project looks like. Edit \filepath{Cargo.toml}:

\begin{minted}{toml}
[package]
name = "blink_bare_metal"
version = "0.1.0"
edition = "2021"

[profile.dev]
opt-level = "z"          # Optimize for size
debug = true

[profile.release]
opt-level = "z"          # Optimize for size
panic = "abort"          # Simplest panic handler

[dependencies]
# We'll add dependencies as needed
\end{minted}

Now modify \filepath{src/main.rs}:

\begin{minted}{rust}
#![no_std]    // Don't link the standard library (no heap, no OS)
#![no_main]   // We define our own entry point

// Prevent Rust from generating main() wrapper
use core::panic::PanicInfo;

/// Our entry point - Cortex-M calls this on reset
#[no_mangle]
pub extern "C" fn main() -> ! {
    // Application logic will go here

    loop {
        // Infinite loop to prevent return
    }
}

/// Simple panic handler - just halt
#[panic_handler]
fn panic(_info: &PanicInfo) -> ! {
    loop {}
}
\end{minted}

\subsection{What is no\_std?}

\begin{notebox}
    \textbf{no\_std} means ``no standard library.'' Our program cannot use:
    \begin{itemize}
        \item Heap allocation (Box, Vec, String)
        \item File I/O
        \item Threads
        \item Networking
    \end{itemize}

    We can still use: \codeinline{core}, slices, static arrays, and direct hardware access.
\end{notebox}

The standard library requires an operating system for features like heap allocation and threads. In bare-metal embedded, we have no OS, so we use only \codeinline{core}, which provides fundamental types and operations that don't require an OS.

\subsection{The Entry Point}

In hosted environments (with an OS), Rust generates a C main() function that calls your main(). In bare metal:

\begin{itemize}
    \item \codeinline{#[no_mangle]} prevents Rust from changing the function name
    \item \codeinline{pub extern "C" fn} ensures the correct calling convention
    \item \codeinline{-> !} means "never returns" (diverging function)
    \item We must provide an infinite loop or halt instruction
\end{itemize}

\subsection{Cross-Compilation}

Our development machine (likely x86\_64) can't execute ARM Cortex-M code. We need to cross-compile:

\begin{minted}{bash}
$ rustup target add thumbv8m.main-none-eabihf
\end{minted}

This installs the Rust standard library for the ARM Cortex-M33 architecture without an OS.

Create \filepath{.cargo/config.toml}:

\begin{minted}{toml}
[build]
target = "thumbv8m.main-none-eabihf"

[target.thumbv8m.main-none-eabihf]
# No `runner` on purpose: `cargo run` would program the chip the moment it builds,
# and on this board that has to go through the guarded script (see the note below).
rustflags = [
  "-C", "link-arg=-Tlink.x",
  "-C", "link-arg=--nmagic",
]

[env]
DEFMT_LOG = "debug"
\end{minted}

\begin{quote}
\textbf{Why this book has no cargo run.} The first 24 KB of flash
belong to the Gecko bootloader. A `runner` in `.cargo/config.toml` turns every
\codeinline{cargo run} into a write, and a single wrong linker script silently relocates
the vector table into that region. This project flashes through
\codeinline{scripts/flash-lnx.sh}: dry run by default, \codeinline{--confirm} to write,
a hard refusal when the image's lowest address is below \codeinline{0x08006000}, and a
read-back verification afterwards. The same script is what re-flashes the board after a
mistake — which matters more than the convenience of \codeinline{cargo run}.
\end{quote}

\subsection{Building the Project}

\begin{minted}{bash}
$ cargo build --release
   Compiling blink_bare_metal v0.1.0
    Finished release [optimized] target(s) in 0.5s
\end{minted}

Check the result:

\begin{minted}{bash}
$ cargo size --release
   text    data     bss     dec     hex filename
     64       0       0      64      40 blink_bare_metal
\end{minted}

Our minimal program is only 64 bytes!

\subsection{Analyzing the Binary}

Use \codeinline{objdump} to examine the machine code:

\begin{minted}{bash}
$ cargo objdump --release -- -d
\end{minted}

This shows the ARM assembly instructions. You'll see something like:

\begin{minted}{text}
bare_metal:     file format elf32-littlearm

Disassembly of section .text:

main:
    b #-4 <main>
    b #-4 <main>
\end{minted}

The two \codeinline{b #-4} instructions are branches to themselves---infinite loops. One is our main() function, the other is the panic handler.

\subsection{What This Demonstrates}

Even at this minimal stage, we've established:

\begin{itemize}
    \item Cross-compilation toolchain setup
    \item Bare-metal runtime configuration
    \item Minimal program that builds successfully
    \item Linker script and startup code (provided by cortex-m-rt)
\end{itemize}

The program won't do anything visible yet (no LEDs blinking), but it represents a working foundation. In the next section, we'll add actual hardware interaction.

\begin{exercise}[Easy]
    Modify the panic handler to increment a counter in an infinite loop. Use a static mut variable and examine how the compiler warns about unsafe code.
\end{exercise}

\begin{exercise}[Medium]
    Research what happens when main() (or the reset handler) returns on a Cortex-M microcontroller. Why is \codeinline{-> !} important?
\end{exercise}

\begin{exercise}[Hard]
    Create a minimal linker script (\filepath{linker.ld}) and build script (\filepath{build.rs}) to understand how the linker allocates memory for your program.
\end{exercise}
\clearpage
